\chapter{Wilson-Line Frobenius and Archimedean Wilson-Surface at Genus \(g\)}
\label{v4-arith:w9-r6:chapter}

\emph{Chapter~\ref{v4-arith:all-genus-koszul} states the conditional
all-genus arithmetic Koszul duality reduction
\(\kappa^{\mathrm{Ar},g}_{\mathsf{G}}+(\kappa^{!})^{\mathrm{Ar},g}_{\mathsf{G}}=0\)
on an arithmetic genus-\(g\) datum
\(\mathcal{D}_{g}=(\mathcal{B}_{g},\mathcal{E}_{g},\varphi_{g})\), with
the genus entering through the rank-\(2g\) Fock grading, the archimedean
\(\Gamma_{\mathbb{C}}(s)^{g}\) factor, and the motivic weight
\(w_{\mathcal{D}_{g}}\), equal to \(1\) in the zeta window, \(2\) in the
curve \(H^{1}\) window, and \(w_{\pi,\mathrm{spin}}\) in the Siegel spin
window. Chapter~\ref{v4-arith:3d-acs-E1bar}
constructs finite-level arithmetic Chern--Simons state sums and states
the profinite quantum ACS hypotheses on
\(\overline{\mathrm{Spec}\,\mathbb{Z}}\), including the Wilson-line
Frobenius axiom, here written in the geometric normalization used by
Euler factors,
\(\langle W_{\rho}(K_{p})\rangle=\mathrm{tr}\,\rho(\mathrm{Frob}_{p}^{\mathrm{geom}})\)
at \(g=0\). Chapter~\ref{v4-arith:w8-a:deninger-chiral-homology}
separates the ordered-incidence candidate from Deninger's conjectural
arithmetic cohomology; their identification and the global Frobenius
\(F_{\infty}\otimes\bigotimes_{p}'\psi_{p}\) remain conditional on the
spectral-flow and boundary-comparison hypotheses. The genus-\(g\)
extension replaces the gauge group
\(GL_{2}(\widehat{\mathbb{Z}})\) by its rank-\(2g\) analog
\(GL_{2g}(\widehat{\mathbb{Z}})\); at each prime \(p\) the defect
operator becomes the genus-\(g\) Wilson line, and the archimedean boundary
operator becomes the genus-\(g\) Wilson-surface \(F^{(g)}_{\infty}\). The genus-\(g\) gauge group is the profinite limit of
\(GL_{2g}(\mathbb{Z}/N\mathbb{Z})\); the rank-\(2g\) local representation
\(\rho_{g}:G_{\mathbb{Q}_{p}}\to\mathrm{GSp}_{2g}\) sources a genus-\(g\)
Wilson loop whose expectation in
\(\mathrm{ACS}^{\mathrm{qu},g}\) is
\(\mathrm{tr}_{\rho_{g}}(\mathrm{Frob}_{p}^{\mathrm{geom}})\); the archimedean
Wilson-surface \(F^{(g)}_{\infty}\) is the rank-\(2g\) Arakelov area-scaling
conjugate to the genus-\(g\) Virasoro zero-mode
\(L^{(g)}_{0}\) with central charge \(g\). The all-genus Deninger
identification is conditional; when its hypotheses are imposed it is
the motivic-weight-\(w_{\mathcal{D}_{g}}\) specialisation of the
character identity of
\Cref{v4-arith:w8-a:thm:all-genus-deninger}. At \(g=0\) the genus
family degenerates to the rank-zero datum. The rank-two Wilson-line
theory of the Riemann framework is an independent ACS sector, not the
value of \(\rho_{g,p}\) at \(g=0\). The compatibility with
Chs.~\ref{v4-arith:w9-r1:wilson-frobenius},
\ref{v4-arith:w9-r2:archimedean-wilson-surface},
\ref{v4-arith:rh-reformulation}, and \ref{v4-arith:tfe-gutzwiller} is
recorded in \Cref{v4-arith:w9-r6:cor:g-zero}.}

\section{Genus-\(g\) Arithmetic Data Recap}
\label{v4-arith:w9-r6:datum}

Fix \(g\geq 0\) and an arithmetic genus-\(g\) datum
\(\mathcal{D}_{g}=(\mathcal{B}_{g},\mathcal{E}_{g},\varphi_{g})\) as in
\Cref{v4-arith:ag:defn:datum}. Four structural features of
\(\mathcal{D}_{g}\) are used throughout.

\begin{enumerate}
\item \emph{Local Galois representation.} At each prime \(p\)
of good reduction for \(\mathcal{D}_{g}\), there is a continuous rank-\(2g\)
representation
\(\rho_{g,p}:G_{\mathbb{Q}_{p}}:=\mathrm{Gal}(\overline{\mathbb{Q}}_{p}/\mathbb{Q}_{p})\to\mathrm{GSp}_{2g}(\overline{\mathbb{Q}}_{\ell})\)
with symplectic Weil pairing; in Window~I, \(\rho_{g,p}=T_{\ell}\mathrm{Jac}(X_{K})_{\mathbb{Q}_{p}}\)
(Tate module of the Jacobian); in Window~III, the prismatic
\(F\)-crystal; in Window~VI, the standard representation is
\((2g+1)\)-dimensional and the spin representation is
\(2^{g}\)-dimensional, so the Wilson label must specify which
representation is used.
\item \emph{Satake parameters \(\{\alpha_{j}^{(p)}\}_{j=1}^{2g}\).}
At an unramified prime, \(\rho_{g,p}(\mathrm{Frob}_{p}^{\mathrm{geom}})\) has
\(2g\) eigenvalues \(\alpha_{1}^{(p)},\dots,\alpha_{2g}^{(p)}\) with
symplectic pairing \(\alpha_{j}^{(p)}\alpha_{2g+1-j}^{(p)}=p^{w_{\mathcal{D}_{g}}-1}\);
these source the standard local Euler factor. In the Siegel spin row,
the spin Euler factor is formed from the spin representation weights,
and the Wilson observable must carry the spin label.
\item \emph{Rank-\(2g\) Fock oscillators per place.} The local factor
\(\mathcal{H}^{(p),g}\) of
\(\mathcal{V}^{\mathrm{prim},g}_{\mathrm{Heis}}=\bigotimes_{p}'\mathcal{H}^{(p),g}\otimes\mathcal{H}^{(\infty),g}\)
is the rank-\(2g\) Fock whose character is the local Euler factor above.
The archimedean factor \(\mathcal{H}^{(\infty),g}\) carries the
rank-\(2g\) archimedean Heisenberg with \(\Gamma\)-factor
\(\Gamma_{\mathbb{C}}(s)^{g}\) at weight \(1\) (Windows~I, V) or the
corresponding \(\Gamma\)-factor of
\Cref{v4-arith:ag:table:rows} at Dedekind or Siegel-spin weight.
\item \emph{Motivic functional-equation centre
\(w_{\mathcal{D}_{g}}\).} The completed \(L\)-function
\(\Lambda_{\mathcal{D}_{g}}(s)\) satisfies
\(\Lambda_{\mathcal{D}_{g}}(s)=\epsilon_{\mathcal{D}_{g}}\Lambda_{\mathcal{D}_{g}}(w_{\mathcal{D}_{g}}-s)\)
by \Cref{v4-arith:ag:thm:all-genus}; purity lies on
\(\mathrm{Re}(s)=w_{\mathcal{D}_{g}}/2\).
\end{enumerate}

\begin{remark}[degeneration at \(g=0\)]
\label{v4-arith:w9-r6:rem:g-zero}
At \(g=0\) the datum \(\mathcal{D}_{0}\) has no non-trivial Galois
representation: \(\mathcal{E}_{0}\) is void, \(\varphi_{0}=\mathrm{id}\),
and there is no rank-\(2g\) local system left to trace. The representation
\(\rho\) in Chapter~\ref{v4-arith:w9-r1:wilson-frobenius} is a separate
\(GL_{2}\)-valued Wilson-line label.
The Heisenberg degenerates to the rank-\(0\) Fock whose local factor is
the trivial Euler line, and the global character \(\xi_{\mathbb{R}}(s)\)
retains only the Bost--Connes primon-gas contribution
\(\zeta(s)\) multiplied by \(\pi^{-s/2}\Gamma(s/2)\) and the Tate-pole
factor \(\tfrac{1}{2}s(s-1)\). The Wilson-line identity of
Chapter~\ref{v4-arith:w9-r1:wilson-frobenius} and the non-trivial
archimedean Wilson-surface of
Chapter~\ref{v4-arith:w9-r2:archimedean-wilson-surface} are compatible
rank-two and rank-one Tate inputs, not consequences of the empty
rank-\(2g\) trace. The comparison is recorded in
Section~\ref{v4-arith:w9-r6:g-zero-recovery}.
\end{remark}

\section{Genus-\(g\) 3D Arithmetic Chern--Simons Theory}
\label{v4-arith:w9-r6:3d-acs-g}

Chapter~\ref{v4-arith:3d-acs-E1bar} builds the 3D arithmetic
Chern--Simons theory on \(\overline{\mathrm{Spec}\,\mathbb{Z}}\) with
gauge group \(GL_{2}(\widehat{\mathbb{Z}})\) in the profinite limit,
level \(k^{\mathrm{Ar}}=\kappa^{\mathrm{Ar}}_{\mathsf{G}}\), boundary WZW
\(\mathcal{V}^{\mathrm{prim}}\), and Wilson-loop axiom
\Cref{v4-arith:3d-acs:conj:qu-acs}(3). The datum \(\mathcal{D}_{g}\) determines
a rank-\(2g\) gauge group \(GL_{2g}(\widehat{\mathbb{Z}})\), a rank-\(2g\)
Koszul level \(k^{\mathrm{Ar},g}\), and a rank-\(2g\) boundary Heisenberg
\(\mathcal{V}^{\mathrm{prim},g}\); the following definition records these data
and inherits the conditional form of the \(g=0\) theory.

\begin{definition}[genus-\(g\) arithmetic Chern--Simons]
\label{v4-arith:w9-r6:def:acs-g}
\ClaimStatusDefinitional. The \emph{genus-\(g\) arithmetic Chern--Simons
theory} is a family of 3D topological gauge theories on an arithmetic
\(3\)-space \(\mathcal{B}_{g}^{(3)}\) parameterised by
\(\mathcal{D}_{g}\), with
\begin{enumerate}
\item \emph{Arithmetic \(3\)-space.} For Window~II at totally imaginary
\(K\) of degree \(2g\):
\(\mathcal{B}_{g}^{(3)}=\overline{\mathrm{Spec}\,\mathcal{O}_{K}}\) with
Artin--Verdier fundamental class
\([\mathcal{B}_{g}^{(3)}]\in H^{3}_{\mathrm{\acute{e}t}}(\mathcal{B}_{g}^{(3)},\widehat{\mathbb{Z}})\)
(Mazur 1973). For Window~I at \(X\to\mathrm{Spec}\,\mathcal{O}_{K}\):
\(\mathcal{B}_{g}^{(3)}=\overline{\mathrm{Spec}\,\mathcal{O}_{K}}\) with
extra fibre cohomology in \(H^{2}\) of the generic fibre collected into
the archimedean boundary via the Arakelov K\"ahler form
\(\mu_{\mathrm{Ar},g}\) of Faltings 1984 (to be defined in
\Cref{v4-arith:w9-r6:def:arakelov-kahler-g}). For Window~VI, Shimura varieties
of \(\mathrm{GSp}_{2g}\) have their own Artin--Verdier-type fundamental
class assembled via \'etale cohomology of the special fibre plus
Deligne 1974 Weil cohomology.
\item \emph{Gauge group.}
\(G_{g,N}:=GL_{2g}(\mathbb{Z}/N\mathbb{Z})\), with profinite limit
\(G_{g,\infty}:=GL_{2g}(\widehat{\mathbb{Z}})\).
\item \emph{Level.}
\(k^{\mathrm{Ar},g}:=\kappa^{\mathrm{Ar},g}_{\mathsf{G}}(\mathcal{D}_{g})\),
the rank-\(2g\) arithmetic Koszul level of
\Cref{v4-arith:ag:thm:all-genus}.
\item \emph{Boundary WZW at the archimedean fibre.}
\(\partial\,\mathrm{ACS}^{\mathrm{qu},g}_{k^{\mathrm{Ar},g}}=\mathcal{V}^{\mathrm{prim},g}\)
in the sense of
\Cref{v4-arith:3d-acs:thm:bulk-boundary-Vprim} scaled to rank \(2g\) per
\Cref{v4-arith:3d-acs:genus-g}.
\end{enumerate}
The finite-level state sums
\(Z^{\mathrm{Ar}}_{G_{g,N},\alpha_{g,N}}\) extend
\Cref{v4-arith:3d-acs:thm:finite-level-ACS} unchanged at each finite
rank; the profinite renormalisation
\(\mathrm{ACS}^{\mathrm{qu},g}_{k^{\mathrm{Ar},g}}\) remains conditional
in the sense of \Cref{v4-arith:3d-acs:conj:qu-acs} extended to rank
\(2g\).
\end{definition}

\begin{remark}[rank scaling and residual profinite limit]
\label{v4-arith:w9-r6:rem:rank-scaling}
The rank-change \(2\to 2g\) does not alter the closed finite-level
state-sum of \Cref{v4-arith:3d-acs:thm:finite-level-ACS}: the field
groupoid
\(\mathcal{F}_{G_{g,N}}(\mathcal{B}_{g}^{(3)})\)
is a finite groupoid for each \(N\), the action
\(S_{\alpha_{g,N}}^{\mathrm{Ar}}\) is the same Artin--Verdier pairing,
and the gluing formula follows from \'etale descent of
rank-\(2g\) local systems. The profinite limit
\(\mathrm{ACS}^{\mathrm{qu},g}\) shares the profinite-renormalisation
residual of \Cref{v4-arith:3d-acs:conj:qu-acs}.
\end{remark}

\begin{proposition}[partition function at genus \(g\)]
\label{v4-arith:w9-r6:prop:partition-g}
\ClaimStatusConditional \emph{(on the genus-\(g\) extension of
\Cref{v4-arith:3d-acs:conj:qu-acs}(2) and the genus-\(g\) BD/CG
comparison of \Cref{v4-arith:3d-acs:cor:BDCG-ordered-comparison})}.
The partition function of
\(\mathrm{ACS}^{\mathrm{qu},g}_{k^{\mathrm{Ar},g}}\) on
\(\mathcal{B}_{g}^{(3)}\) equals the completed \(L\)-function
\begin{equation}
\label{v4-arith:w9-r6:eq:partition-g}
Z_{\mathrm{ACS}^{\mathrm{qu},g}}\bigl(\mathcal{B}_{g}^{(3)};k^{\mathrm{Ar},g}\bigr)
\;=\;\Lambda_{\mathcal{D}_{g}}\bigl(k^{\mathrm{Ar},g}\bigr),
\end{equation}
with \(\Lambda_{\mathcal{D}_{g}}\) the row-specific completed \(L\)-function
of \Cref{v4-arith:ag:table:rows} (Dedekind \(\xi_{K}\), Hasse--Weil
\(\Lambda(s,X_{K})\), or Siegel spin \(\Lambda(s,\pi,\mathrm{spin})\)).
\end{proposition}

\begin{proof}
The rank-\(2g\) ordered incidence bar \(B^{\mathrm{ord},g}_{\bullet}\) on
\(\mathbb{P}\cup\{\infty\}\) with per-place factor \(\mathcal{H}^{(p),g}\)
is the rank-\(2g\) generalisation of the bar of
\Cref{v4-arith:3d-acs:defn:E1-bar}. The local character at \(p\) is
\(\chi_{p}^{(g),\mathrm{std}}(s)=\prod_{j=1}^{2g}(1-\alpha_{j}^{(p)}p^{-s})^{-1}\)
for the standard Wilson observable. In the Siegel spin row the local
character is instead the spin Euler factor
\(\chi_{p}^{(g),\mathrm{spin}}(s)=\prod_{\epsilon}(1-\beta_{\epsilon}^{(p)}p^{-s})^{-1}\)
over the weights of the spin representation. The archimedean character
is the row-specific \(\Gamma\)-factor. The genus-\(g\) extension of
\Cref{v4-arith:tfe:lem:ord-bar} gives
\(H_{\bullet}(B^{\mathrm{ord},g})=\Lambda_{\mathcal{D}_{g}}(s)\). Under
the conditional genus-\(g\) BD/CG comparison and the renormalised
profinite limit, the bulk partition function coincides with this
homology evaluated at \(s=k^{\mathrm{Ar},g}\).
\end{proof}

\section{Wilson-Line Frobenius Identity at Genus \(g\)}
\label{v4-arith:w9-r6:wilson-line-g}

The gauge group \(G_{g,\infty}=GL_{2g}(\widehat{\mathbb{Z}})\) of
\(\mathrm{ACS}^{\mathrm{qu},g}\) contains the prime knots \(K_{p}\) as
defect insertions. Each \(K_{p}\) carries a holonomy valued in the
rank-\(2g\) local Galois representation \(\rho_{g,p}\); the question is
whether the expectation of the associated defect operator recovers
geometric Frobenius \(\mathrm{Frob}_{p}^{\mathrm{geom}}\). Arithmetic
Frobenius is its inverse in this convention. The answer at rank \(2g\) is the
following identity.

\begin{definition}[genus-\(g\) Wilson line]
\label{v4-arith:w9-r6:def:wilson-line-g}
\ClaimStatusDefinitional. Fix a prime \(p\) of good reduction for
\(\mathcal{D}_{g}\); let \(K_{p}\subset\mathcal{B}_{g}^{(3)}\) denote the
prime-knot associated to \(p\) under the Morishita dictionary. For a
continuous representation
\(\rho_{g}:G_{\mathbb{Q}_{p}}\to\mathrm{GSp}_{2g}(\overline{\mathbb{Q}}_{\ell})\)
of rank \(2g\), the \emph{genus-\(g\) Wilson line} along \(K_{p}\) is the
defect operator
\begin{equation}
\label{v4-arith:w9-r6:eq:wilson-line-def}
W^{(g)}_{\rho_{g}}(K_{p})\;:=\;\mathrm{tr}_{\rho_{g}}\bigl(\mathrm{Hol}_{K_{p}}(\mathsf{A})\bigr),
\end{equation}
with \(\mathsf{A}\in\mathcal{F}_{G_{g,\infty}}(\mathcal{B}_{g}^{(3)})\)
the bulk rank-\(2g\) gauge field in the profinite limit and
\(\mathrm{Hol}_{K_{p}}\) the path-ordered holonomy around the meridian
\(K_{p}\), whose finite local factors are the rank-\(2g\)
Bost--Connes Adams operations \(\psi^{(g)}_{p}\).
\end{definition}

\begin{theorem}[genus-\(g\) Wilson-line Frobenius identity]
\label{v4-arith:w9-r6:thm:wilson-frobenius-g}
\ClaimStatusProvedHereConditional \emph{(on the rank-\(2g\) extension of
\Cref{v4-arith:3d-acs:conj:qu-acs}(3) and on rank-\(2g\) local Langlands
at unramified \(p\) of good reduction for \(\mathcal{D}_{g}\), per the
genus-\(g\) Fargues--Scholze arXiv:2102.13459 input
\Cref{v4-arith:ag:residual}(5))}. For every unramified prime \(p\) of
good reduction for \(\mathcal{D}_{g}\),
\begin{equation}
\label{v4-arith:w9-r6:eq:wilson-frobenius-g}
\langle W^{(g)}_{\rho_{g}}(K_{p})\rangle_{\mathrm{ACS}^{\mathrm{qu},g}}
\;=\;\mathrm{tr}_{\rho_{g}}\bigl(\mathrm{Frob}_{p}^{\mathrm{geom}}\bigr)
\;=\;\sum_{j=1}^{2g}\alpha_{j}^{(p)}.
\end{equation}
Equivalently, the Wilson-line defect expectation equals the character value
\(\chi_{\rho_{g}}(\mathrm{Frob}_{p}^{\mathrm{geom}})\) in the standard representation.
The Siegel spin \(L\)-factor is obtained from the spin Wilson observable,
not from this rank-\(2g\) standard trace.
\end{theorem}

\begin{proof}
Three steps.

\emph{Step 1: Morishita-dictionary reduction.} The prime-knot
\(K_{p}\subset\mathcal{B}_{g}^{(3)}\) corresponds, under the
Mazur--Morishita dictionary (Mazur 1973 ENS~6; Morishita 2012
arXiv:0904.3399), to the inertia-generated meridian around the closed
point \(\mathfrak{p}\) above \(p\). Restriction of the rank-\(2g\) gauge
field \(\mathsf{A}\) to a tubular neighbourhood of \(K_{p}\) is determined
by the local Galois representation \(\rho_{g,p}\) via the rank-\(2g\)
generalisation of \Cref{v4-arith:3d-acs:cor:finite-wilson}: at finite
level \(G_{g,N}\), the defect expectation equals
\(\chi_{\rho_{g}\,\mathrm{mod}\,N}(\mathrm{Frob}_{p}^{\mathrm{geom}})\) for the induced
mod-\(N\) representation.

\emph{Step 2: Profinite assembly.} The compatible tower
\(\{\chi_{\rho_{g}\,\mathrm{mod}\,N}(\mathrm{Frob}_{p}^{\mathrm{geom}})\}_{N}\) assembles
under the renormalised profinite limit of
\Cref{v4-arith:3d-acs:conj:qu-acs} to the continuous Frobenius trace
\(\mathrm{tr}_{\rho_{g}}(\mathrm{Frob}_{p}^{\mathrm{geom}})\). This requires the
rank-\(2g\) refinement of the Wilson-line axiom of
\Cref{v4-arith:3d-acs:conj:qu-acs}(3), which is the residual
hypothesis; it is the same assumption that supplies
\Cref{v4-arith:w8-a:lem:wilson-finite} at \(g=0\).

\emph{Step 3: Bost--Connes Adams identification.} Borger 2009
arXiv:0906.3146 \S8 identifies each \(\psi_{p}\) with geometric
Frobenius on the \(\lambda\)-ring envelope of the Tate motive; the
rank-\(2g\) generalisation replaces the scalar action on
\(\mathcal{H}_{p}^{\mathrm{BC}}\) by the \(2g\)-fold generating-series
action with Satake parameters
\(\{\alpha_{j}^{(p)}\}_{j=1}^{2g}\). Under the unramified hypothesis,
\(\mathrm{tr}\,\rho_{g}(\mathrm{Frob}_{p}^{\mathrm{geom}})=\sum_{j=1}^{2g}\alpha_{j}^{(p)}\)
by definition of the Satake parameters. Combining with Step~1--2 gives
\eqref{v4-arith:w9-r6:eq:wilson-frobenius-g}.
\end{proof}

\begin{remark}[independence of sources in the Wilson-line Frobenius identity]
\label{v4-arith:w9-r6:rem:independence-wilson-line}
The three inputs of \Cref{v4-arith:w9-r6:thm:wilson-frobenius-g} are independent.
The rank-\(2g\) Bost--Connes Adams action of Borger 2009 \S8 derives from absolute
\(\lambda\)-ring theory; the rank-\(2g\) local Langlands of Fargues--Scholze 2021
arXiv:2102.13459 from \(p\)-adic geometric Langlands; the Morishita knot-prime
dictionary from Mazur 1973 plus Morishita 2012 arithmetic topology. The three
source families are pairwise disjoint.

The Satake parameter normalisation is: unramified \(\rho_{g,p}\) has eigenvalues
\(\alpha_{j}^{(p)}\) with symplectic pairing
\(\alpha_{j}^{(p)}\alpha_{2g+1-j}^{(p)}=p^{w_{\mathcal{D}_{g}}-1}\); the
Frobenius trace is \(\sum_{j=1}^{2g}\alpha_{j}^{(p)}\), the character value of
\(\rho_{g}\) itself at \(\mathrm{Frob}_{p}^{\mathrm{geom}}\), not the character of
\(\bigwedge^{k}\rho_{g}\) for any \(k\geq 2\).

The Wilson-line defect at finite level lies in
\(\mathrm{Sh}_{\tau_{\Sigma}(\alpha_{g,N})}(\mathcal{F}_{G_{g,N}}(\Sigma))\)
for \(\Sigma=\partial(\text{tubular neighbourhood of }K_{p})\) per
\Cref{v4-arith:3d-acs:thm:finite-bulk-boundary}; the rank-\(2g\) setting is
structurally identical. The Wilson line is functorial under tensor product of
representations at the classical level:
\(W^{(g)}_{\rho_{1}\otimes\rho_{2}}=W^{(g)}_{\rho_{1}}\cdot W^{(g)}_{\rho_{2}}\);
at the quantum level this uses the factorisation algebra structure of
\(\mathcal{V}^{\mathrm{prim},g}\) (rank-\(2g\) Heisenberg), consistent with
Ch.~\ref{v4-arith:all-genus-koszul} Window~V.

The residual hypotheses are: (a) profinite renormalisation of
\(\{Z^{\mathrm{Ar}}_{G_{g,N},\alpha_{g,N}}\}_{N}\) and (b) compatible tower
\(\{\alpha_{g,N}\}_{N}\), inherited from \Cref{v4-arith:3d-acs:residual}(2); and
(c) rank-\(2g\) local Langlands at \(p\), i.e.\ \Cref{v4-arith:ag:residual}(5).
Identification of the bulk holonomy \(\mathrm{Hol}_{K_{p}}(\mathsf{A})\) with
\(\rho_{g,p}\) passes through the Artin--Verdier fundamental-class pairing of
\Cref{v4-arith:3d-acs:def:finite-field-groupoid}, unchanged at rank \(2g\);
the good-reduction hypothesis excludes wild ramification and singular fibres.
\end{remark}

\section{Archimedean Wilson-Surface at Genus \(g\)}
\label{v4-arith:w9-r6:wilson-surface-g}

At the archimedean place, a three-dimensional ACS boundary is a real
surface. For the curve row, the 3-space \(\mathcal{B}_{g}^{(3)}\) has a
closed-surface boundary \(\Sigma_{\infty}=\mathcal{B}_{g}^{(3)}\mid_{\infty}\)
of real dimension \(2\), with genus \(g\) entering through
\(H^{1}(\Sigma_{\infty})\) of rank \(2g\), rather than through the
dimension of the boundary. A defect operator on \(\Sigma_{\infty}\) must integrate an area form
rather than a holonomy along a loop: the Faltings 1984 arithmetic K\"ahler form
\(\mu_{\mathrm{Ar},g}\) on the archimedean fibre is the natural input, since it
encodes the metric data of \(\mathcal{D}_{g}\) at the archimedean place and
degenerates to zero at \(g=0\), matching the trivial archimedean boundary of the
Riemann \(\xi\)-function framework.

\begin{definition}[genus-\(g\) archimedean Arakelov K\"ahler form]
\label{v4-arith:w9-r6:def:arakelov-kahler-g}
\ClaimStatusDefinitional. The \emph{genus-\(g\) Arakelov K\"ahler form}
on the archimedean fibre \(\mathcal{B}_{g}^{(3)}\mid_{\infty}\) is
\begin{equation}
\label{v4-arith:w9-r6:eq:arakelov-kahler-g}
\mu_{\mathrm{Ar},g}\;:=\;\frac{i}{2g}\sum_{j=1}^{g}\omega_{j}\wedge\overline{\omega_{j}},
\end{equation}
with \(\{\omega_{j}\}_{j=1}^{g}\) a basis of holomorphic differentials on
the archimedean fibre
\(X_{\infty}(\mathbb{C})\hookrightarrow\mathcal{B}_{g}^{(3)}\mid_{\infty}\)
(for Window~I) or \(g\) copies of the archimedean Haar form on the
Minkowski torus \(T_{K}\) (for Window~II) or the natural Riemannian
volume form scaled per Siegel-modular signature (for Window~VI). This
is the Faltings 1984 arithmetic K\"ahler form (Faltings 1984
\emph{Ann.~Math.}~119) on the generic fibre; at \(g=0\) the sum is
empty and
\(\mu_{\mathrm{Ar},0}=0\).
\end{definition}

\begin{definition}[archimedean Wilson-surface at genus \(g\)]
\label{v4-arith:w9-r6:def:wilson-surface-g}
\ClaimStatusDefinitional. The \emph{genus-\(g\) archimedean
Wilson-surface} is the bulk area-scaling defect
\begin{equation}
\label{v4-arith:w9-r6:eq:wilson-surface-g-def}
F^{(g)}_{\infty}\;:=\;\exp\Bigl(-k^{\mathrm{Ar},g}\int_{\Sigma_{\infty}}\mu_{\mathrm{Ar},g}\Bigr),
\end{equation}
where \(\Sigma_{\infty}\subset\mathcal{B}_{g}^{(3)}\) is the archimedean
fibre \(\mathcal{B}_{g}^{(3)}\mid_{\infty}\) viewed as a closed
\(2\)-surface with the Faltings K\"ahler form of
\Cref{v4-arith:w9-r6:def:arakelov-kahler-g}. In operator-theoretic
form on
\(\mathcal{H}^{(\infty),g}\),
\begin{equation}
\label{v4-arith:w9-r6:eq:wilson-surface-g-op}
F^{(g)}_{\infty}\;=\;\exp\bigl(-k^{\mathrm{Ar},g}\sum_{j=1}^{g}\widetilde L_{0,j}^{\mathrm{Ar}}\bigr),
\end{equation}
for \(\widetilde L_{0,j}^{\mathrm{Ar}}\) the centred Fock scaling
zero-mode on the \(j\)-th archimedean Heisenberg factor
\(\mathcal{H}^{(\infty),g}_{j}\). The uncentred scaling generator has
Mellin eigenvalue \(s_{j}\); the tilde subtracts the critical-line
centre \(w_{\mathcal{D}_{g}}/2\).
\end{definition}

\begin{theorem}[conjugacy to genus-\(g\) Virasoro]
\label{v4-arith:w9-r6:thm:virasoro-conjugacy-g}
\ClaimStatusProvedHereConditional \emph{(on the rank-\(g\) extension of
the primary-Virasoro reconciliation
\Cref{v4-arith:rh:conj:primary-Virasoro} per
\Cref{v4-arith:ag:residual}(7), and on the rank-\(g\) Arakelov--Mellin
unitarity of \Cref{v4-arith:w9-r6:lem:mellin-unitary-g} below)}. The
genus-\(g\) archimedean Wilson-surface \(F^{(g)}_{\infty}\) is conjugate
to the genus-\(g\) arithmetic Virasoro zero-mode
\(L^{(g)}_{0}\) with central charge \(c^{\mathrm{Ar},g}=g\):
\begin{equation}
\label{v4-arith:w9-r6:eq:virasoro-conjugacy-g}
F^{(g)}_{\infty}\;=\;U^{(g)}\,\exp\bigl(-L^{(g)}_{0}\bigr)\,(U^{(g)})^{-1},
\end{equation}
with \(U^{(g)}\) the rank-\(g\) Mellin-transform unitary of
\Cref{v4-arith:w9-r6:lem:mellin-unitary-g}. Equivalently, the Sugawara
construction on the rank-\(2g\) archimedean Heisenberg produces
\(L^{(g)}_{0}\) as the arithmetic Virasoro zero-mode sourced by the
rank-\(2g\) archimedean Fock; the Arakelov area-scaling
\(F^{(g)}_{\infty}\) is its Mellin conjugate.
\end{theorem}

\begin{lemma}[rank-\(g\) Arakelov--Mellin unitarity]
\label{v4-arith:w9-r6:lem:mellin-unitary-g}
\ClaimStatusProvedHereConditional \emph{(on the rank-\(g\) Weil--Petersson
isometry lift of Ch.~\ref{v4-arith:w5-a:chapter} restricted to the
archimedean fibre; genus-\(g\) Petersson inner-product of
\Cref{v4-arith:ag:residual}(2))}. There exists a unitary
\(U^{(g)}:L^{2}(\mathbb{R}_{+}^{g},\prod_{j=1}^{g}d^{\times}x_{j})\to L^{2}(\mathrm{Re}\,s_{j}=w_{\mathcal{D}_{g}}/2,\,j=1,\dots,g)\)
defined by the \(g\)-fold tensor of Mellin transforms
\begin{equation}
\label{v4-arith:w9-r6:eq:mellin-g}
U^{(g)}f(s_{1},\dots,s_{g})\;=\;\int_{\mathbb{R}_{+}^{g}}
f(x_{1},\dots,x_{g})\prod_{j=1}^{g}x_{j}^{s_{j}}\,
\prod_{j=1}^{g}d^{\times}x_{j},
\end{equation}
conjugating the centred rank-\(g\) Arakelov area-scaling flow
\(\partial^{\circ}_{t_{j}}=\partial/\partial\log\|\cdot\|_{\infty,j}
-w_{\mathcal{D}_{g}}/2\) to
multiplication by \(s_{j}-w_{\mathcal{D}_{g}}/2\) on the critical
poly-line.
\end{lemma}

\begin{proof}
The Faltings K\"ahler form
\eqref{v4-arith:w9-r6:eq:arakelov-kahler-g} decomposes as a sum of
\(g\) orthogonal rank-\(1\) pieces under the Hodge decomposition
\(\mathcal{E}_{g}\otimes\mathbb{C}=\mathcal{E}_{g}^{1,0}\oplus\mathcal{E}_{g}^{0,1}\);
each piece carries a rank-\(1\) Arakelov scaling flow
\(t_{j}\mapsto e^{-t_{j}}\|\cdot\|_{\infty,j}\) with Mellin variable
\(s_{j}\). The rank-\(1\) unitary is
\(\mathrm{Mellin}:L^{2}(\mathbb{R}_{+},d^{\times}x)\to L^{2}(\mathrm{Re}\,s=w_{\mathcal{D}_{g}}/2)\),
classical (Tate 1950; Titchmarsh 1986 Chapter~14); the \(g\)-fold tensor
is unitary by Fubini on the product measure. The scaling flow is
diagonalised by Mellin at the centre
\(\mathrm{Re}\,s_{j}=w_{\mathcal{D}_{g}}/2\) to match
\Cref{v4-arith:w9-r6:datum}(4).
\end{proof}

\begin{proof}[Proof of \Cref{v4-arith:w9-r6:thm:virasoro-conjugacy-g}]
Under \(U^{(g)}\), the archimedean Wilson-surface
\eqref{v4-arith:w9-r6:eq:wilson-surface-g-op} becomes
\(\exp(-k^{\mathrm{Ar},g}\sum_{j=1}^{g}(s_{j}-w_{\mathcal{D}_{g}}/2))\) on
the product Mellin space; this coincides with
\(\exp(-L^{(g)}_{0})\) where
\(L^{(g)}_{0}=\sum_{j=1}^{g}\widetilde L_{0,j}^{\mathrm{Ar}}\) is the rank-\(g\)
arithmetic Virasoro zero-mode on the product archimedean Heisenberg
Fock. The Sugawara construction of
\Cref{v4-arith:3d-acs:thm:km-lift} scales rank-linearly per
Ch.~\ref{v4-arith:all-genus-koszul} (central charge \(c^{\mathrm{Ar},g}=g\));
the rank-\(g\) primary sector admits a Virasoro zero-mode which is the
natural lift of \(\widetilde L_{0,j}^{\mathrm{Ar}}\) on each archimedean Heisenberg
factor. Summing over factors gives the rank-\(g\) zero-mode
\(L^{(g)}_{0}\); this is the \Cref{v4-arith:rh:conj:primary-Virasoro}
reconciliation at rank \(g\).
\end{proof}

\begin{remark}[Wilson-surface realisation as bulk \(\partial_{k}\)-derivative]
\label{v4-arith:w9-r6:rem:wilson-surface-realisation}
In the 3D ACS framework, the archimedean Wilson-surface is the bulk
partition-function derivative along the arithmetic-level deformation
restricted to the archimedean fibre:
\(F^{(g)}_{\infty}=\partial Z_{\mathrm{ACS}^{\mathrm{qu},g}}/\partial k^{\mathrm{Ar},g}\mid_{\mathrm{archimedean}}\),
by the area-generation scheme of
\Cref{v4-arith:w8-a:rem:wilson-infinity} scaled to rank \(2g\). This is
the arithmetic analog of the Witten-Verlinde area operator on complex
Riemann surfaces.
\end{remark}

\begin{remark}[independence of sources and sign verification for the archimedean Wilson-surface]
\label{v4-arith:w9-r6:rem:independence-wilson-surface}
The three primary inputs are independent: Faltings 1984 Arakelov
K\"ahler form (Arakelov geometry), \Cref{v4-arith:rh:conj:primary-Virasoro}
(primary Virasoro sector, chiral algebra methods), and Tate 1950 Mellin
unitarity (classical analysis). The sign in
\(F^{(g)}_{\infty}=\exp(-k^{\mathrm{Ar},g}\int\mu_{\mathrm{Ar},g})\) matches
the Mellin convention \(\partial_{t}\leftrightarrow -s+w_{\mathcal{D}_{g}}/2\)
on the critical line; the Faltings normalisation
\eqref{v4-arith:w9-r6:eq:arakelov-kahler-g} with prefactor \((2g)^{-1}\)
ensures the total area integrates to \(g\) times the area per factor.

The Wilson-surface lives in bulk codimension-\(1\) defects of
\(\mathrm{ACS}^{\mathrm{qu},g}\); at finite level in
\(\mathrm{Sh}_{\tau(\alpha_{g,N})}(\mathcal{F}_{G_{g,N}}(\Sigma_{\infty}))\)
restricted to the archimedean surface. The Wilson-surface is the tensor product
of \(g\) rank-\(1\) Wilson-surfaces when the Faltings K\"ahler form splits as a
direct sum: this holds for Jacobians (Window~I) and Minkowski tori (Window~II)
by Hodge decomposition, but Window~VI requires the Siegel--Maass metric on
\(\mathcal{A}_{g}\) (residual 2 of this chapter).

The residual hypotheses are: rank-\(g\) primary-Virasoro reconciliation
(\Cref{v4-arith:ag:residual}(7)); rank-\(g\) Weil--Petersson isometry lift
(\Cref{v4-arith:ag:residual}(2)); rank-\(g\) Mellin unitarity
(\Cref{v4-arith:w9-r6:lem:mellin-unitary-g}), which is conditioned on the
Weil--Petersson lift.
\end{remark}

\section{Global Frobenius at Genus \(g\) and the Deninger Identification}
\label{v4-arith:w9-r6:deninger-g}

The Wilson-line and Wilson-surface operators at each place are local.
To obtain the Deninger character identity
\(\chi(\mathcal{V}^{\mathrm{prim},g})(s)=\Lambda_{\mathcal{D}_{g}}(s)\), one
must assemble them into a single global operator on the conditional
Deninger cohomology \(H^{\bullet}_{\Delta,g}\). The restricted tensor
product over all places is therefore a formal global Frobenius until
the descent and spectral-flow hypotheses are imposed.

\begin{definition}[global genus-\(g\) Frobenius]
\label{v4-arith:w9-r6:def:global-F-g}
\ClaimStatusDefinitional. The \emph{formal global genus-\(g\)
Frobenius} is the restricted-tensor symbol
\begin{equation}
\label{v4-arith:w9-r6:eq:global-F-g}
F^{(g)}\;:=\;F^{(g)}_{\infty}\otimes\bigotimes_{p}{}^{\prime}\,\psi^{(g)}_{p},
\end{equation}
acting first on
\(\mathcal{V}^{\mathrm{prim},g}_{\mathrm{Heis}}=\mathcal{H}^{(\infty),g}\otimes\bigotimes_{p}'\mathcal{H}^{(p),g}\).
Its descent to a Deninger cohomology
\(H^{\bullet}_{\Delta,g}\) is not part of the definition; it is exactly
one of the conditional inputs in
\Cref{v4-arith:w8-a:thm:all-genus-deninger}.
\end{definition}

\begin{theorem}[genus-\(g\) Wilson assembly of the global Frobenius]
\label{v4-arith:w9-r6:thm:global-F-g-wilson}
\ClaimStatusProvedHereConditional \emph{(on
\Cref{v4-arith:w9-r6:thm:wilson-frobenius-g},
\Cref{v4-arith:w9-r6:thm:virasoro-conjugacy-g}, and the
Tate-restricted-tensor compatibility F2.c at rank \(2g\), per
\Cref{v4-arith:cl:F2-closure} scaled)}. The formal global Frobenius
\(F^{(g)}\) admits bulk Wilson realisations:
\begin{enumerate}
\item \emph{Finite places.} \(\psi^{(g)}_{p}\) is realised as the
genus-\(g\) Wilson-line observable along the knot \(K_{p}\):
\(\mathrm{tr}_{\rho_{g}}(\psi^{(g)}_{p})=\langle W^{(g)}_{\rho_{g}}(K_{p})\rangle_{\mathrm{ACS}^{\mathrm{qu},g}}
=\mathrm{tr}_{\rho_{g}}(\mathrm{Frob}_{p}^{\mathrm{geom}})\).
\item \emph{Archimedean place.} \(F^{(g)}_{\infty}\) is realised as the
archimedean Wilson-surface of
\Cref{v4-arith:w9-r6:def:wilson-surface-g}, conjugate to
\(\exp(-L^{(g)}_{0})\) with central charge \(c^{\mathrm{Ar},g}=g\).
\end{enumerate}
\end{theorem}

\begin{proof}
Item~(1) follows from \Cref{v4-arith:w9-r6:thm:wilson-frobenius-g}
applied per-prime, assembled under Tate restricted-tensor
compatibility (F2.c) across almost-all primes. At genus \(g\), F2.c
scales to rank \(2g\) without structural change: the spherical unit
\(|0\rangle_{p}\) of the rank-\(2g\) local Heisenberg is a unique
\(\psi^{(g)}_{p}\)-eigenvector up to normalisation \(p^{-w_{\mathcal{D}_{g}}/2}\).

Item~(2) is \Cref{v4-arith:w9-r6:thm:virasoro-conjugacy-g}.

The conditional global descent to \(H^{\bullet}_{\Delta,g}\) uses the rank-\(2g\) bar
complex of \Cref{v4-arith:w8-a:def:deninger-g}, invariant under each
\(\psi^{(g)}_{p}\) by multiplicativity and under \(F^{(g)}_{\infty}\) by the
graded-Virasoro action.
\end{proof}

\begin{theorem}[all-genus Deninger identification at motivic weight \(w_{\mathcal{D}_{g}}\)]
\label{v4-arith:w9-r6:thm:deninger-g-full}
\ClaimStatusProvedHereConditional \emph{(on
\Cref{v4-arith:w8-a:thm:all-genus-deninger}(1)--(4),
\Cref{v4-arith:w9-r6:thm:wilson-frobenius-g},
\Cref{v4-arith:w9-r6:thm:virasoro-conjugacy-g}, and
\Cref{v4-arith:w9-r6:prop:partition-g})}. For every \(g\geq 0\) and every
arithmetic genus-\(g\) datum
\(\mathcal{D}_{g}=(\mathcal{B}_{g},\mathcal{E}_{g},\varphi_{g})\):
\begin{enumerate}
\item \emph{Triple labelling.} The Vol~IV genus-\(g\) Deninger
cohomology \(H^{\bullet}_{\Delta,g}\), the boundary chiral homology
of \(\mathrm{ACS}^{\mathrm{qu},g}_{k^{\mathrm{Ar},g}}\) on
\(\mathcal{B}_{g}^{(3)}\), and the rank-\(2g\) arithmetic chiral homology
\(H^{\bullet}_{\mathrm{ch}}(\mathcal{B}_{g}^{(3)},\mathsf{H}^{\mathrm{Ar},g})\)
(with \(\mathsf{H}^{\mathrm{Ar},g}\) the arithmetic Heisenberg coefficient
sheaf of \Cref{v4-arith:w8-a:def:deninger-g})
coincide on the primary rank-\(2g\) Heisenberg domain.
\item \emph{Character identity with motivic weight.}
\begin{equation}
\label{v4-arith:w9-r6:eq:character-g}
\chi\bigl(\mathcal{V}^{\mathrm{prim},g}\bigr)(s)
\;=\;\det\nolimits_{\infty}\bigl(s\,\mathrm{id}-\Theta^{(g)}\mid H^{\bullet}_{\Delta,g}\bigr)^{(-1)^{\bullet}}
\;=\;\Lambda_{\mathcal{D}_{g}}(s),
\end{equation}
where \(\Lambda_{\mathcal{D}_{g}}\) has motivic weight
\(w_{\mathcal{D}_{g}}\) per \Cref{v4-arith:ag:table:rows}; in the
Siegel spin row this is \(w_{\pi,\mathrm{spin}}\), not a universal
function of \(g\).
\item \emph{Wilson-observable form of the Frobenius.} The local
factors of \(F^{(g)}\) appearing in
\eqref{v4-arith:w9-r6:eq:character-g} admit the Wilson-line /
Wilson-surface realisation of
\Cref{v4-arith:w9-r6:thm:global-F-g-wilson}.
\item \emph{Purity at \(\mathrm{Re}(s)=w_{\mathcal{D}_{g}}/2\).}
The genus-\(g\) purity axiom of \Cref{v4-arith:w8-a:deninger-wishlist}(6)
becomes \(F^{(g)}\)-weight \(w_{\mathcal{D}_{g}}/2\) on
\(H^{1}_{\Delta,g}\); equivalently the motivic Riemann Hypothesis for
the row \(\mathcal{D}_{g}\) (\Cref{v4-arith:ag:cor:all-genus-FRH}).
\end{enumerate}
\end{theorem}

\begin{proof}
Item~(1) is \Cref{v4-arith:w8-a:thm:all-genus-deninger}(1) unchanged,
scaled to rank \(2g\).

Item~(2): the regularised determinant \(\det_{\infty}(s\,\mathrm{id}-\Theta^{(g)}\mid
H^{\bullet}_{\Delta,g})^{(-1)^{\bullet}}\) decomposes across Fock degrees:
\(H^{0}_{\Delta,g}\cong\mathbb{C}_{0}\) contributes \(s\);
\(H^{2}_{\Delta,g}\cong\mathbb{C}_{w_{\mathcal{D}_{g}}}\) (Tate
twist by the motivic weight) contributes \(s-w_{\mathcal{D}_{g}}\);
\(H^{1}_{\Delta,g}\) carries the
regularised product \(\prod_{\rho}(s-\rho)\) over the non-trivial
zeros of \(\Lambda_{\mathcal{D}_{g}}\). Multiplying and normalising by
the \(\pi^{-gs/2}\Gamma_{\mathbb{C}}(s)^{g}\) absorbed in the definition
of \(F^{(g)}_{\infty}\) gives \(\Lambda_{\mathcal{D}_{g}}(s)\) per
\Cref{v4-arith:w8-a:thm:all-genus-deninger}(3).

Item~(3) is \Cref{v4-arith:w9-r6:thm:global-F-g-wilson}.

Item~(4) is the rank-\(g\) analog of
\Cref{v4-arith:w8-a:thm:purity-equivalence} via
\Cref{v4-arith:ag:cor:all-genus-FRH}: weight-\(w_{\mathcal{D}_{g}}/2\)
purity is the statement that all non-trivial zeros of
\(\Lambda_{\mathcal{D}_{g}}\) lie on
\(\mathrm{Re}(s)=w_{\mathcal{D}_{g}}/2\), subsuming the classical RH at
\(g=0\), Hasse--Weil at weight \(2\), and Siegel-spin at
\(w_{\pi,\mathrm{spin}}\).
\end{proof}

\begin{remark}[source independence and residuals for the all-genus Deninger identification]
\label{v4-arith:w9-r6:rem:independence-deninger}
The five primary inputs are pairwise independent: Deninger 1998
(foliated dynamical systems), Bost--Connes 1995 (KMS partition on the
id\`ele class group), Fargues--Scholze 2021 (rank-\(2g\) local Langlands),
Faltings 1984 (Arakelov K\"ahler form), Arthur 2013 (endoscopic classification
for Siegel spin). The motivic-weight normalisation is: Dedekind \(w=1\),
Hasse--Weil \(w=2\), Siegel spin \(w_{\pi,\mathrm{spin}}\), per
\Cref{v4-arith:ag:table:rows};
the alternating-sign Euler product \(\det(\cdot)^{(-1)^{\bullet}}\) follows
Deninger's determinant convention; the central charge \(g\) of the Sugawara
construction matches the archimedean rank.

The rank-\(2g\) bar \(B^{\mathrm{Ar},g}(\mathcal{V}^{\mathrm{prim},g})\) lives
in graded \(\mathbb{C}\)-vector spaces with rank-\(2g\) Heisenberg Fock grading;
the \(E_{1}\)-ordered bar on \(\mathbb{P}\cup\{\infty\}\) is rank-enhanced per
\Cref{v4-arith:w9-r6:prop:partition-g}. Functoriality under pullback along
\(\overline{\mathrm{Spec}\,\mathcal{O}_{K}}\to\overline{\mathrm{Spec}\,\mathbb{Z}}\)
requires \(g=\tfrac{1}{2}[K:\mathbb{Q}]\) at totally imaginary \(K\)
(Window~II); for Shimura (Window~VI), functoriality under Hecke correspondences
holds because \(F^{(g)}\) commutes with Hecke operators on the rank-\(2g\)
automorphic side.

The full residual list consists of all items of \Cref{v4-arith:ag:residual}
(rank-\(g\) Positselski, rank-\(g\) Koszul--Petersson, \(GL_{2g}\) local
geometric Langlands, rank-\(g\) Virasoro bootstrap, rank-\(g\) F2.c) and
all items of \Cref{v4-arith:3d-acs:residual} (profinite ACS renormalisation,
boundary comparison, BD/CG comparison, positivity), each scaled to rank \(2g\).
The identification of the \(H^{1}_{\Delta,g}\) spectrum with the non-trivial
zeros of \(\Lambda_{\mathcal{D}_{g}}\) requires rank-\(g\) H-P\(^{\mathrm{Ar}}\):
a self-adjoint lift of \(L^{(g)}_{0}\) to a Hilbert completion of
\(H^{1}_{\Delta,g}\) in the rank-\(g\) arithmetic Petersson inner product
(\Cref{v4-arith:rh:hyp:HPAr} at rank \(g\)); this is open.
\end{remark}

\section{Uniform \(g=0\) Compatibility}
\label{v4-arith:w9-r6:g-zero-recovery}

At \(g=0\), the constructions of Sections~\ref{v4-arith:w9-r6:3d-acs-g}--\ref{v4-arith:w9-r6:deninger-g}
degenerate to the rank-zero skeleton inside the Riemann framework of
Chs.~\ref{v4-arith:3d-acs-E1bar} and
\ref{v4-arith:w8-a:deninger-chiral-homology}. The non-trivial
rank-two Wilson sectors are compared, not derived, below.

\begin{itemize}
\item \emph{Archimedean Wilson-surface.} At \(g=0\),
\(\mu_{\mathrm{Ar},0}=0\) (empty Hodge sum); the Wilson-surface
\eqref{v4-arith:w9-r6:eq:wilson-surface-g-def} degenerates to the
trivial defect
\(\exp(-k^{\mathrm{Ar}}\cdot 0)=1\). The non-trivial Arakelov scaling of
\Cref{v4-arith:w8-a:def:F-infinity} and
\Cref{v4-arith:w9-r2:def:F-infty-WS} belongs to the rank-one Tate
Mellin sector; it is not obtained by substituting \(g=0\) in the
rank-\(g\) surface operator.
\item \emph{Wilson-line Frobenius.} At \(g=0\), \(\rho_{0}\) is the
empty rank-\(0\) datum, so the genus-\(g\) Wilson-line identity has no
finite local trace. The identity
\(\langle W_{\rho}(K_{p})\rangle=\mathrm{tr}\,\rho(\mathrm{Frob}_{p}^{\mathrm{geom}})\)
of \Cref{v4-arith:3d-acs:conj:qu-acs}(3) and
\Cref{v4-arith:w9-r1:thm:main} uses an independent continuous
\(GL_{2}\)-label \(\rho\); it is compatible with, but not a
specialisation of, \eqref{v4-arith:w9-r6:eq:wilson-frobenius-g}.
\item \emph{All-genus Deninger identification.} At \(g=0\),
\(\Lambda_{\mathcal{D}_{0}}(s)=\Lambda_{\zeta}(s)\), and the
pole-cleared determinant is
\(\xi_{\mathbb R}(s)=\tfrac12s(s-1)\Lambda_{\zeta}(s)\). The motivic
weight is \(w_{\mathcal{D}_{0}}=1\); the character identity
\eqref{v4-arith:w9-r6:eq:character-g} collapses to the \(g=0\)
identity \eqref{v4-arith:w8-a:eq:character} of
\Cref{v4-arith:w8-a:thm:character-identity}; the purity locus is
\(\mathrm{Re}(s)=1/2\), the Riemann critical line.
\end{itemize}

\begin{corollary}[uniform \(g=0\) compatibility]
\label{v4-arith:w9-r6:cor:g-zero}
\ClaimStatusProvedHere. The genus-\(g\) constructions of
Sections~\ref{v4-arith:w9-r6:3d-acs-g},
\ref{v4-arith:w9-r6:wilson-line-g},
\ref{v4-arith:w9-r6:wilson-surface-g} and
\ref{v4-arith:w9-r6:deninger-g} degenerate at \(g=0\) to the
rank-zero Tate skeleton. The rank-two constructions of
Chs.~\ref{v4-arith:3d-acs-E1bar},
\ref{v4-arith:w8-a:deninger-chiral-homology},
\ref{v4-arith:w9-r1:wilson-frobenius}, and
\ref{v4-arith:w9-r2:archimedean-wilson-surface} are separate compatible
inputs:
Wilson-line Frobenius \Cref{v4-arith:3d-acs:conj:qu-acs}(3) and
\Cref{v4-arith:w9-r1:thm:main}; archimedean Wilson-surface
\Cref{v4-arith:w8-a:def:F-infinity} and
\Cref{v4-arith:w9-r2:def:F-infty-WS} with conjugacy
\Cref{v4-arith:w9-r2:thm:conjugacy}; Deninger character identity
\Cref{v4-arith:w8-a:thm:character-identity}. No rank-two Wilson-line
label is produced by the rank-zero degeneration.
\end{corollary}

\begin{proof}
Each item in the bulleted list above separates the rank-zero
degeneration from the independent rank-two ACS input. No analytic
correction is required for the degeneration: the rank-\(2g\) Fock is
empty, the Hodge decomposition \(\mathcal{E}_{0}\otimes\mathbb{C}=0\),
and the motivic weight is \(w_{\mathcal{D}_{0}}=1\).
\end{proof}

\begin{remark}[Hasse--Weil and Siegel-spin RH at genus \(g\geq 1\)]
\label{v4-arith:w9-r6:rem:hw-ss-rh}
At \(g\geq 1\), the Deninger-chiral purity locus
\(\mathrm{Re}(s)=w_{\mathcal{D}_{g}}/2\) is the motivic Riemann
Hypothesis for the row:
\begin{itemize}
\item \(w_{\mathcal{D}_{g}}=1\) (Dedekind, Window~II at totally imaginary
\(K\) of degree \(2g\)): the generalised Riemann Hypothesis for
\(\zeta_{K}(s)\).
\item \(w_{\mathcal{D}_{g}}=2\) (Hasse--Weil, Windows~I, V at the Jacobian
of a curve of genus \(g\)): the Hasse--Weil Riemann Hypothesis on the
weight-\(1\) motive \(H^{1}(\mathrm{Jac}(X_{K}))\), critical line
\(\mathrm{Re}(s)=1\).
\item \(w_{\mathcal{D}_{g}}=w_{\pi,\mathrm{spin}}\) (Siegel spin, Window~VI at cusp form
\(\pi\) on \(\mathrm{GSp}_{2g}\)): the Siegel-spin Riemann Hypothesis,
critical line \(\mathrm{Re}(s)=w_{\pi,\mathrm{spin}}/2\).
\end{itemize}
The Wilson-line / Wilson-surface realisation converts each case to a
bulk topological-gauge-theoretic spectral statement on the
partition-function zeros of \(\mathrm{ACS}^{\mathrm{qu},g}_{k^{\mathrm{Ar},g}}\)
on the corresponding arithmetic 3-space \(\mathcal{B}_{g}^{(3)}\).
\end{remark}

\section{Residual Conditions}
\label{v4-arith:w9-r6:residual}

The genus-\(g\) extension inherits residuals from
Ch.~\ref{v4-arith:all-genus-koszul} and Ch.~\ref{v4-arith:3d-acs-E1bar}
scaled to rank \(2g\), and four residuals are specific to this chapter.

\begin{enumerate}
\item \emph{Rank-\(2g\) local Langlands compatibility at ramified
primes.} \Cref{v4-arith:w9-r6:thm:wilson-frobenius-g} holds at
unramified primes of good reduction for \(\mathcal{D}_{g}\). At ramified
primes, the identity requires a local twist term depending on the
Artin conductor and the Fargues--Scholze / Genestier--Lafforgue
correspondence for \(\mathrm{GSp}_{2g}\). It is a non-formal theorem,
extending \Cref{v4-arith:ag:residual}(5).
\item \emph{Rank-\(g\) Faltings K\"ahler splitting for Shimura Window~VI.}
The archimedean Wilson-surface
\eqref{v4-arith:w9-r6:eq:wilson-surface-g-def} presumes the Faltings
K\"ahler form splits as a direct sum of \(g\) orthogonal rank-\(1\) pieces;
this holds for Jacobians (Window~I) and Minkowski tori (Window~II) by
Hodge decomposition, but requires the Siegel--Maass metric on
\(\mathcal{A}_{g}\) (Kudla 1997 \emph{Ann.~Math.}~146) for Window~VI.
It is a non-formal theorem.
\item \emph{Rank-\(g\) Bost--Connes Adams compatibility.} Borger 2009
\S8's identification of \(\psi_{p}\) with geometric Frobenius on the
Tate \(\lambda\)-motive scales to rank \(2g\) via the \(\lambda\)-ring structure on
the rank-\(2g\) arithmetic Heisenberg. Formally immediate at each
finite prime; the restricted-tensor compatibility across almost all
primes is F2.c at rank \(2g\), inheriting
\Cref{v4-arith:ag:residual}(6). It is a non-formal theorem.
\item \emph{Rank-\(g\) H-P\(^{\mathrm{Ar}}\).} Self-adjoint lift of
\(L^{(g)}_{0}\) to a Hilbert completion of
\(H^{1}_{\Delta,g}\) in the rank-\(g\) arithmetic Petersson inner product.
This is the direct rank-\(g\) analog of
\Cref{v4-arith:rh:hyp:HPAr}. Domain: inherited from H-P\(^{\mathrm{Ar}}\)
plus \Cref{v4-arith:ag:residual}(2)--(7). The residual is the motivic
Riemann Hypothesis in Hilbert--P\'olya form.
\end{enumerate}

\begin{remark}[uniform rank structure]
\label{v4-arith:w9-r6:rem:uniform-rank}
The genus enters through the rank of the gauge group, the archimedean Fock
rank, and the motivic weight; the Koszul self-duality identity
\(\kappa^{\mathrm{Ar},g}+(\kappa^{!})^{\mathrm{Ar},g}=0\) is fixed on
the residual domain of \Cref{v4-arith:ag:thm:all-genus}; the bulk
partition function is the row-specific completed \(L\)-function
\(\Lambda_{\mathcal{D}_{g}}\); and the motivic Riemann Hypothesis for
the row becomes the topological spectral condition
\(\mathrm{Re}(z)=w_{\mathcal{D}_{g}}/2\) on the non-trivial zeros
\(z\) of \(\Lambda_{\mathcal{D}_{g}}\), i.e.\ the conditional
\(F^{(g)}\)-weight purity on \(H^{1}_{\Delta,g}\)
(\Cref{v4-arith:w9-r6:thm:deninger-g-full}(4)).
\end{remark}

\section{Primary Sources for the Genus-\(g\) Extension}
\label{v4-arith:w9-r6:sources}

Seven primary sources underpin the genus-\(g\) extension, each disjoint
from the others in the sense of \Cref{v4-arith:ag:defn:datum}.

\begin{itemize}
\item \emph{Faltings 1984} \emph{Calculus on arithmetic surfaces},
\emph{Ann.~Math.}~119, 387--424: Arakelov K\"ahler form
\(\mu_{\mathrm{Ar},g}\) of
\Cref{v4-arith:w9-r6:def:arakelov-kahler-g}; Faltings--Hriljac
arithmetic Hodge index; Faltings \(\delta\)-invariant.
\item \emph{Fargues--Scholze 2021} arXiv:2102.13459: rank-\(2g\) local
Langlands for \(GL_{2g}/\mathbb{Q}_{p}\); Zhu 2020 arXiv:2008.02998
and Genestier--Lafforgue 2017 arXiv:1709.00978 extend to
\(\mathrm{GSp}_{2g}\).
\item \emph{Bhatt--Scholze 2022} \emph{Ann.~Math.}~196: prismatic
Frobenius on rank-\(2g\) \(F\)-crystals;
\Cref{v4-arith:ag:w3-fargues-scholze} Window~III input.
\item \emph{Borger 2009} arXiv:0906.3146 \S8: Bost--Connes Adams
\(\psi_{p}\) as geometric Frobenius on the Tate \(\lambda\)-motive; rank-\(2g\) scaling via the
\(\lambda\)-ring structure on the Tate motive.
\item \emph{Arthur 2013} \emph{The Endoscopic Classification}: Siegel
spin \(L\)-functions and functional equations for
\(\mathrm{GSp}_{2g}\)-automorphic cuspidals.
\item \emph{Kudla 1997} \emph{Ann.~Math.}~146; \emph{Kudla--Rapoport
2011}: Siegel--Maass metric on \(\mathcal{A}_{g}\) and
rank-\(g\) Faltings splitting for Window~VI.
\item \emph{Morishita 2012} \emph{Knots and Primes}
arXiv:0904.3399; \emph{Mazur 1973} ENS~6: knot-prime dictionary for
the Wilson-line identification at finite primes; Artin--Verdier
fundamental class for the arithmetic 3-space.
\end{itemize}

Pairwise disjointness of these sources, as in the independent-source
list of \Cref{v4-arith:ag:orthogonal}, is preserved at rank \(2g\)
by the same inheritance argument: the rank-change is a scaling
of the per-place Heisenberg oscillators, not a source substitution.

The residuals of \Cref{v4-arith:w9-r6:residual} together with those of
Chs.~\ref{v4-arith:all-genus-koszul}, \ref{v4-arith:3d-acs-E1bar},
and \ref{v4-arith:w8-a:deninger-chiral-homology} constitute the full
condition list for a theorem-grade proof of the genus-\(g\) Wilson-line,
Wilson-surface, and Deninger identification. Four conditions remain open:
(a) rank-\(2g\) profinite ACS renormalisation, (b) rank-\(2g\) BD/CG
comparison, (c) rank-\(g\) Weil--Petersson isometry, (d) rank-\(g\)
H-P\(^{\mathrm{Ar}}\). Each is recorded as a conditional input in its
natural chapter of origin.
